Đúng hướng rồi. Với mức trừu tượng bạn muốn, tôi sẽ tổ chức phần formalization thành một chuỗi rất tự nhiên: **khái niệm nguồn → khái niệm đích → hai phép chuyển → tính đúng đắn**. Cách này cũng gần với phong cách bài bạn gửi: trước hết định nghĩa các mô hình bằng tuple và điều kiện well-formedness, sau đó mới đưa phương pháp và tiêu chí conformance/correctness.  

## 1. Definition: Class Metamodel, Class Model, and Object Model

Mở đầu bằng một định nghĩa rất ngắn:

$$
\boxed{
MM_C=\langle \mathcal C,\mathcal A,\mathcal R,\mathcal H,\mathcal O\rangle
}
$$

Trong đó \(MM_C\) chỉ quy định các **loại phần tử có thể xuất hiện trong một class model**: class, attribute, association, inheritance và OCL constraint.

Sau đó không cần định nghĩa lại một metamodel thứ hai. Chỉ nói:

> A class model \(M\) is an instance of \(MM_C\).

Ký hiệu:

$$
\boxed{
M\models MM_C
}
$$

và có thể biểu diễn:

$$
\boxed{
M=\langle C_M,A_M,R_M,H_M,O(M)\rangle.
}
$$

Ở đây:

$$
O(M)=\{o_1,\ldots,o_n\}
$$

là tập OCL invariants thuộc \(M\).

Sau đó giới thiệu object model/snapshot:

$$
\boxed{
S=\langle Obj_S,Slot_S,Link_S\rangle
}
$$

và:

$$
\boxed{
S\models M.
}
$$

Nên giải thích ngắn:

* \(Obj_S\): objects;
* \(Slot_S\): attribute values;
* \(Link_S\): association instances.

Điểm cần nói rõ là:

$$
S\models M
$$

chỉ là **structural conformance**, không có nghĩa tất cả OCL trong \(M\) đều đúng.

Đây là chỗ quan trọng để reader hiểu tại sao vẫn có thể tồn tại:

$$
Viol(O,S)\neq\varnothing.
$$

---

## 2. Well-formedness và OCL phù hợp với \(M\)

Không cần trình bày hàng chục typing rules.

Chỉ cần một judgement tổng quát:

$$
\boxed{
M;\Gamma\vdash e:\tau
}
$$

nghĩa là biểu thức OCL \(e\):

* dùng đúng phần tử khai báo trong \(M\);
* variable/binder đúng scope;
* attribute/navigation resolve được;
* toán tử dùng đúng type;
* multiplicity/navigation phù hợp.

Với invariant:

$$
o=\langle K,e\rangle
$$

ta định nghĩa:

$$
\boxed{
WF_{OCL}(M,o)
\iff
K\in C_M
\land
M;\{self:K\}\vdash e:Boolean.
}
\tag{1}
$$

Sau đó có thể nói:

$$
\boxed{
WF(M)
\iff
M\models MM_C
\land
\forall o\in O(M).\ WF_{OCL}(M,o).
}
\tag{2}
$$

Đây là đủ. Toàn bộ proof chi tiết của typing/name resolution để Isabelle artifact.

---

# 3. Definition: Graph Metamodel and Graph

Phía graph giữ đúng kiến trúc hai tầng.

Định nghĩa:

$$
\boxed{
MM_G
}
$$

là metamodel của property graph, quy định một graph gồm node, relationship, labels, types và properties.

Một graph cụ thể:

$$
\boxed{
G=
\langle
V,E,\lambda,\rho,src,tgt,\pi
\rangle
}
$$

trong đó:

$$
V
$$

là nodes,

$$
E
$$

là relationships,

$$
\lambda:V\rightarrow 2^{Label}
$$

là labels,

$$
\rho:E\rightarrow RelType
$$

là relationship type,

$$
src,tgt:E\rightarrow V
$$

là endpoints,

và:

$$
\pi:(V\cup E)\times Key\rightharpoonup Value
$$

là property function.

Graph hợp lệ khi:

$$
\boxed{
G\models MM_G.
}
\tag{3}
$$

Không cần giải thích `node_wf`, `edge_wf`, finiteness... ở main text.

Một câu prose là đủ:

> \(G\models MM_G\) requires finite well-formed nodes and relationships, valid relationship endpoints, labels, types, and property values.

---

# 4. Một query hợp lệ là gì?

Query không cần một metamodel riêng.

Chỉ cần ký hiệu:

$$
\boxed{
WF_Q(M,\Sigma,q)
}
$$

nghĩa là \(q\):

* đúng cú pháp query language;
* chỉ sử dụng labels, relationship types và property keys được ánh xạ từ \(M\) thông qua \(\Sigma\);
* navigation/direction phù hợp với representation specification;
* kết quả trả về đúng dạng violation nodes/IDs.

Nếu muốn ngắn hơn nữa:

$$
\boxed{
q\in \mathcal Q_{\Sigma}(M)
}
\tag{4}
$$

với:

$$
\mathcal Q_{\Sigma}(M)
$$

là tập các query well-formed đối với \(M\) và representation \(\Sigma\).

Tôi thích notation này hơn vì paper sẽ nhẹ hơn.

---

# 5. Representation specification

Trước hai phép chuyển, giới thiệu một lần:

$$
\boxed{
\Sigma=
\langle
\Sigma_C,\Sigma_A,\Sigma_R,\Sigma_D,\Sigma_{id}
\rangle
}
$$

Trong đó:

* \(\Sigma_C\): class → label;
* \(\Sigma_A\): attribute → property;
* \(\Sigma_R\): association → relationship type;
* \(\Sigma_D\): logical/physical direction;
* \(\Sigma_{id}\): object identity encoding.

Sau này cả \(T_G\) và \(T_Q\) đều dùng cùng \(\Sigma\).

Đây là cầu nối giữa hai transformation.

---

# 6. Transformation 1: Snapshot-to-Graph

Không gọi đây là Definition.

Có thể đặt subsection:

### A. Snapshot-to-Graph Transformation

Giới thiệu:

$$
\boxed{
T_G^\Sigma:(M,S)\mapsto G
}
\tag{5}
$$

với:

$$
M\models MM_C,
\qquad
S\models M.
$$

Transformation tạo:

$$
\boxed{
G=T_G^\Sigma(M,S)
}
$$

và yêu cầu:

$$
\boxed{
G\models MM_G.
}
\tag{6}
$$

Sau đó dùng một bảng nhỏ.

| Source element              | Target graph element    | Mapping rule                  |
| --------------------------- | ----------------------- | ----------------------------- |
| Object \(x\in Obj_S\)       | Node \(v_x\in V\)       | \(x\mapsto v_x\)              |
| Class of \(x\)              | Node label              | \(c\mapsto\Sigma_C(c)\)       |
| Attribute value \(a(x)\)    | Node property           | \(a\mapsto\Sigma_A(a)\)       |
| Association link \(r(x,y)\) | Relationship            | \(r\mapsto\Sigma_R(r)\)       |
| Link direction              | Physical edge direction | determined by \(\Sigma_D(r)\) |
| Object identifier           | Stable graph identifier | determined by \(\Sigma_{id}\) |

Chỉ cần vậy.

Sau bảng có thể tổng kết bằng:

$$
\boxed{
\gamma_\Sigma:
Obj_S\rightarrow V_G
}
$$

là object-to-node correspondence induced by \(T_G^\Sigma\).

Nếu builder tạo đúng một node cho mỗi object:

$$
\boxed{
\gamma_\Sigma
\text{ is bijective}.
}
\tag{7}
$$

---

# 7. Transformation 2: OCL-to-Query

Subsection riêng:

### B. OCL-to-Query Transformation

Giới thiệu:

$$
\boxed{
T_Q^\Sigma:(M,o)\mapsto q,
\qquad o\in O(M).
}
\tag{8}
$$

Transformation chỉ được áp dụng nếu:

$$
WF_{OCL}(M,o).
$$

Và kết quả phải thỏa:

$$
\boxed{
q\in\mathcal Q_\Sigma(M).
}
\tag{9}
$$

Sau đó thêm bảng transformation rules.

| OCL construct       | Query construct        | Ý nghĩa                                       |
| ------------------- | ---------------------- | --------------------------------------------- |
| Context class \(K\) | node scan              | scan nodes labelled \(\Sigma_C(K)\)           |
| Attribute \(e.a\)   | property access        | access \(\Sigma_A(a)\)                        |
| Navigation \(e.r\)  | relationship traversal | use \(\Sigma_R(r)\) and \(\Sigma_D(r)\)       |
| Boolean operators   | graph predicate        | preserve Boolean semantics                    |
| Comparison          | comparison predicate   | preserve typed comparison                     |
| `select`            | filtering              | retain matching elements                      |
| `collect`           | projection             | project expression values                     |
| `forAll`            | universal predicate    | preserve OCL iterator semantics               |
| `exists`            | existential predicate  | preserve OCL iterator semantics               |
| Invariant \(o\)     | violation query        | return context nodes for which \(o\neq true\) |

Nếu muốn nhấn mạnh implementation có các intermediate forms nhưng không muốn làm nặng main text, chỉ viết một câu:

$$
T_Q^\Sigma
=
Realize_\Sigma\circ Lower_\Sigma\circ Front_M.
$$

Không cần khai triển Core và Q ở đây.

---

# 8. Evaluation functions

Đến đây tôi nghĩ **không cần gọi là Definition riêng**.

Chỉ giới thiệu hai hàm quan sát.

OCL:

$$
Eval_O(o,S)
$$

trả về tập object vi phạm:

$$
\boxed{
Eval_O(o,S)\subseteq Obj_S.
}
\tag{10}
$$

Graph query:

$$
Eval_Q(q,G)
$$

trả về tập graph nodes được query xác định là vi phạm:

$$
\boxed{
Eval_Q(q,G)\subseteq V_G.
}
\tag{11}
$$

Một câu prose là đủ:

> \(Eval_O\) denotes standard OCL evaluation over a snapshot, while \(Eval_Q\) denotes execution of the generated query over the transformed graph.

Không cần formalize semantics chi tiết ở đây.

---

# 9. Bây giờ mới Definition: Correctness

Đây mới nên là Definition quan trọng của section.

### Definition X (Transformation Correctness)

Cho:

$$
M\models MM_C,
$$

$$
S\models M,
$$

$$
o\in O(M),
$$

và:

$$
WF_{OCL}(M,o).
$$

Đặt:

$$
G=T_G^\Sigma(M,S),
$$

$$
q=T_Q^\Sigma(M,o).
$$

Phương pháp được gọi là **correct** nếu tồn tại song ánh:

$$
\gamma_\Sigma:
Obj_S
\overset{\sim}{\rightarrow}
V_G
$$

sao cho:

$$
\boxed{
\gamma_\Sigma
\left[
Eval_O(o,S)
\right]
=
Eval_Q(q,G).
}
\tag{12}
$$

Đây là formula chính.

---

## 10. Dạng pointwise còn dễ hiểu hơn

Ngay sau (12), thêm:

$$
\boxed{
\forall x\in Obj_S:
\quad
x\in Eval_O(o,S)
\iff
\gamma_\Sigma(x)\in Eval_Q(q,G).
}
\tag{13}
$$

Rồi giải thích bằng lời:

> An object violates an OCL invariant in the source snapshot if and only if its corresponding graph node is returned by the generated query.

Câu này reviewer đọc là hiểu ngay.

---

# 11. Cuối cùng là theorem

Sau Definition Correctness mới đưa theorem:

### Theorem X (Correctness of the Transformation)

$$
\boxed{
\begin{aligned}
&M\models MM_C
\land
S\models M
\land
WF(M)
\\
&\Longrightarrow
\\
&G=T_G^\Sigma(M,S)
\land
G\models MM_G
\\
&\land\
\forall o\in O(M):
\\
&
\gamma_\Sigma
[
Eval_O(o,S)
]
=
Eval_Q
(
T_Q^\Sigma(M,o),
G
).
\end{aligned}
}
\tag{14}
$$

Đây là theorem lớn duy nhất main paper cần.

---

# Tôi sẽ tổ chức section của paper như thế này

```text
4. Formalization of the Transformation

4.1 Source Modeling Domain
    Definition 1. Class Metamodel and Class Model
    Definition 2. Object Model / Snapshot
    Well-formedness and OCL admissibility

4.2 Target Graph Domain
    Definition 3. Graph Metamodel and Property Graph
    Well-formed graph and well-formed query

4.3 Representation Specification
    Σ

4.4 Snapshot-to-Graph Transformation
    TGΣ : (M,S) → G
    Table I. Snapshot-to-Graph mapping rules

4.5 OCL-to-Query Transformation
    TQΣ : (M,o) → q
    Table II. OCL-to-Query mapping rules

4.6 Semantic Observations
    EvalO(o,S)
    EvalQ(q,G)

4.7 Correctness
    Definition 4. Transformation Correctness
    Theorem 1. Exact Preservation of Violations
```

Tôi nghĩ đây là cấu trúc **vừa đủ formal, vừa đủ trừu tượng, và rất dễ đọc**.

Điểm quan trọng nhất là: **Definition chỉ dành cho những “khái niệm” thực sự cần reader nhớ** — `Class Model`, `Snapshot`, `Graph`, `Correctness`. Hai transformation \(T_G,T_Q\) là **phương pháp**, nên trình bày bằng hàm + bảng mapping rules, đúng như bạn yêu cầu.
